% !TEX program = pdflatex
% ======================================================================
%  Curriculum Vitae -- Mats Petter Kahl
%
%  Build locally:  latexmk -pdf CV_Mats_Kahl.tex   (or pdflatex twice)
%
%  On GitHub, every push that changes this file rebuilds the PDF and
%  publishes it on GitHub Pages (.github/workflows/build-cv.yml), so the
%  link in \latestcv always serves the current version.
% ======================================================================
\documentclass[11pt,a4paper]{article}

% ---- Personal settings -------------------------------------------------
\newcommand{\cvname}{Mats Petter Kahl}
\newcommand{\latestcv}{https://matskahl.github.io/cv/CV_Mats_Kahl.pdf}
\newcommand{\latestcvshort}{matskahl.github.io/cv}

% ---- Fonts, page, colours ----------------------------------------------
\usepackage[T1]{fontenc}
\usepackage[utf8]{inputenc}
\usepackage{XCharter}
\usepackage[protrusion=true,expansion=true,tracking=true]{microtype}
\SetTracking{encoding=*,shape=sc}{40}
\usepackage[a4paper,top=1.9cm,bottom=2.4cm,left=2.2cm,right=2.2cm,footskip=1.1cm]{geometry}
\usepackage{xcolor}
\definecolor{accent}{HTML}{1F3A68}
\definecolor{muted}{HTML}{5F6670}
\usepackage{array}
\usepackage{ragged2e}
\usepackage{enumitem}
\usepackage{titlesec}
\usepackage{fancyhdr}
\usepackage{lastpage}
\usepackage[hyphens]{url}
\usepackage[colorlinks=true,linkcolor=accent,urlcolor=accent]{hyperref}
\hypersetup{
  pdftitle={Curriculum Vitae -- \cvname},
  pdfauthor={\cvname},
  pdfsubject={Curriculum Vitae},
  pdfkeywords={Industrial Organization, Energy Economics, Public Economics, Empirical Microeconomics}
}
\urlstyle{same}
\setlength{\parindent}{0pt}
\setlength{\JustifyingParindent}{0pt}

% ---- Date of this version (set at compile time) -----------------------
\newcommand{\cvmonth}{\ifcase\month\or January\or February\or March\or April\or May\or June\or July\or August\or September\or October\or November\or December\fi}
\newcommand{\cvdate}{\the\day~\cvmonth\ \the\year}

% ---- Section headings --------------------------------------------------
\setcounter{secnumdepth}{0}
\titleformat{\section}
  {\color{accent}\large\scshape}{}{0pt}{}
  [\vspace{-0.35ex}{\color{accent!35}\titlerule[0.6pt]}]
\titlespacing*{\section}{0pt}{2.6ex plus 0.6ex minus 0.3ex}{1.3ex}

% ---- Lists: dates in a left column, content aligned to its right -------
% Lists are set ragged right (which avoids hyphenation); abstracts are justified.
\newlength{\datewidth}
\AtBeginDocument{\settowidth{\datewidth}{2025--present\hspace{1.4em}}}

\newlist{entries}{itemize}{1}
\setlist[entries]{label={}, leftmargin=\datewidth, labelwidth=\datewidth,
  labelsep=0pt, align=left, topsep=0.2ex, partopsep=0pt, parsep=0pt,
  itemsep=0.55ex, before=\raggedright}
\newcommand{\entry}[2]{\item[#1]#2}

\newlist{papers}{itemize}{1}
\setlist[papers]{label={}, leftmargin=\datewidth, topsep=0.2ex,
  partopsep=0pt, parsep=0pt, itemsep=1.1ex, before=\raggedright}

% \paper{url}{title}{coauthors or empty}{venue}
\newcommand{\paper}[4]{\item ``\href{#1}{#2}''\if\relax\detokenize{#3}\relax\else\ (with~#3)\fi\\
  #4}
% \wip{title}{coauthors or empty}
\newcommand{\wip}[2]{\item ``#1''\if\relax\detokenize{#2}\relax\else\ (with~#2)\fi}

\newcommand{\abstracttext}[1]{\par\vspace{0.5ex}{\small\justifying #1\par}}
\newcommand{\note}[1]{\par\vspace{0.3ex}{\small\color{muted}#1\par}}

% \referee{name}{title}{address lines}{email}
\newcommand{\referee}[4]{\begin{minipage}[t]{0.48\linewidth}\raggedright
  \textbf{#1}\\#2\\#3\\\href{mailto:#4}{#4}\end{minipage}}

% ---- Footer ------------------------------------------------------------
\pagestyle{fancy}
\fancyhf{}
\renewcommand{\headrulewidth}{0pt}
\fancyfoot[L]{\footnotesize\color{muted}\cvname\ \textperiodcentered\ Curriculum Vitae\ \textperiodcentered\ \cvdate}
\fancyfoot[R]{\footnotesize\color{muted}Latest version: \href{\latestcv}{\latestcvshort}\quad\thepage\,/\,\pageref*{LastPage}}

% ======================================================================
\begin{document}

% ---- Header ------------------------------------------------------------
{\color{accent}\fontsize{25}{29}\selectfont\scshape \cvname}\hfill
{\color{muted}\scshape Curriculum Vitae}

\vspace{0.4ex}
{\color{accent}\rule{\textwidth}{0.8pt}}

\vspace{0.6ex}
\begin{minipage}[t]{0.47\textwidth}
  \textbf{Visiting Professor}\\
  Department of Economics\\
  Universidad Carlos III de Madrid\\
  C/ Madrid 126, Office 11.2.03\\
  28903 Getafe (Madrid), Spain
\end{minipage}\hfill
\begin{minipage}[t]{0.51\textwidth}
  \raggedleft
  \begin{tabular}[t]{@{}>{\color{muted}}r@{\hspace{0.9em}}l@{}}
    Email       & \href{mailto:mkahl@eco.uc3m.es}{mkahl@eco.uc3m.es}\\
                & \href{mailto:mats.p.kahl@gmail.com}{mats.p.kahl@gmail.com}\\
    Website     & \href{https://sites.google.com/view/mats-kahl}{sites.google.com/view/mats-kahl}\\
    Profiles    & \href{https://orcid.org/0000-0003-3387-6386}{ORCID}\ \textperiodcentered\
                  \href{https://scholar.google.com/citations?user=SzT-tKQAAAAJ}{Google Scholar}\\
    Citizenship & German
  \end{tabular}
\end{minipage}

\vspace{1.2ex}
{\small\color{muted}This version: \cvdate. The most recent version of this CV is always available at
\href{\latestcv}{\latestcvshort}.}

% ---- Research fields ---------------------------------------------------
\section{Research Fields}
\begin{entries}
  \entry{}{Industrial Organization, Energy Economics, Public Economics, Empirical~Microeconomics}
\end{entries}

% ---- Positions ---------------------------------------------------------
\section{Academic Positions}
\begin{entries}
  \entry{2025--present}{\textbf{Visiting Professor}, Department of Economics,
    Universidad Carlos III de Madrid}
  \entry{2025}{\textbf{Visiting Researcher}, Universidad Carlos III de Madrid,
    March--April (host: Philipp~Denter)}
  \entry{2019--2025}{\textbf{Postdoctoral Researcher} (2025) and \textbf{Doctoral Researcher}
    (2019--2025), Institute of Economics, Leuphana University Lüneburg}
\end{entries}

% ---- Education ---------------------------------------------------------
\section{Education}
\begin{entries}
  \entry{2025}{\textbf{Ph.D.\ in Economics}, Leuphana University Lüneburg, \textit{summa cum laude}\\
    Dissertation: \href{https://doi.org/10.48548/pubdata-1659}{\textit{Competition in the German Market for Gasoline Retailing}}\\
    Advisors: Thomas Wein, Boris Hirsch}
  \entry{2019}{\textbf{M.A.\ in Public Economics, Law and Politics}, Leuphana University Lüneburg,
    with~distinction\\
    Exchange semester at the University of Lucerne, Switzerland, 2017}
  \entry{2016}{\textbf{B.Sc.\ in Economics}, University of Mannheim}
\end{entries}

% ---- Job market papers -------------------------------------------------
\section{Job Market Papers}
\begin{papers}
  \item \href{https://github.com/MatsKahl/charging/releases/latest/download/Kahl_Public_Entry_Announced.pdf}{\textbf{Public Entry, Announced: How State-Procured Infrastructure Affects Private Investment}}
  \abstracttext{In 2021, the German government announced where it would place about 1,000 subsidized
    fast-charging sites. This announcement came at least two years before the first of them opened.
    Would private investors retreat, race to pre-empt, or follow the public entrant? Combining the
    government's planning map with 29 archived versions of the federal charging register, I find that
    the announcement did not deter private investment. At motorway rest areas, a loss of more than
    0.22 private charging points within 5\,km can be ruled out. Once public sites opened, public
    capacity was largely additional, raising total local capacity by 0.74 points per public charging
    point (0.90 after adjusting for pre-opening trends). If some public sites would have been built
    privately anyway, true additionality would have been lower. A parsimonious model sheds light on
    the pattern: when the pace of building, rather than demand, limits private investment, public
    capacity is largely additional.}
  \note{Presented at: Applied Reading Group, Universidad Carlos III de Madrid; Research Seminar,
    Universidad Autónoma de Madrid}

  \item \href{https://raw.githubusercontent.com/MatsKahl/spillovers/main/Kahl_Cross_Border_Competition.pdf}{\textbf{Cross-Border Competition and Tax Pass-Through: Evidence from Germany's 2022 Fuel Tax Cut}}
  \abstracttext{How does a national tax cut pass through at an international border on the treated side? How
    does it spill over to the untreated side? I study Germany's 2022 fuel-tax cut using station-level price
    data and exact driving distances to French competitors. German diesel stations within 10\,km of a French
    gas station cut prices by 0.96 cents per liter (95\% confidence interval: 0.43--1.49) more than stations
    further than 50\,km inland. That corresponds to a pass-through premium of about 6 percentage points
    (95\% confidence interval: 2.6--8.9). The premium appears at the start of the cut, is largest within a
    few kilometers of the border, and disappears after the tax returns. French stations near the border also
    cut prices by a similar amount, but they were already falling relative to inland France before the cut.
    The premium arises where the tax cut closed the cross-border price gap, that is, for diesel, and it does
    not appear where the gap stayed wide, that is, for gasoline. A model with treated domestic and untreated
    foreign rivals shows that, at a given intensity of cross-border substitution, feedback through foreign
    prices cannot account for the premium on its own. The evidence is consistent with a change in that
    intensity: cross-border competition that becomes active when the price gap is small. The incidence of a
    national tax cut varies within a country with local competitive conditions.}
  \note{Presented at: EEA Annual Congress 2026 (Dublin); Applied Reading Group, Universidad Carlos III
    de Madrid. Accepted at: SAEe Simposio 2026 (Las Palmas de Gran Canaria)}
\end{papers}

% ---- Publications ------------------------------------------------------
\section{Publications}
\begin{papers}
  \paper{https://doi.org/10.1111/joie.12412}
    {How to Reach the Paradise? Inside the Edgeworth Cycle and Why a Gasoline Station Is the First to Raise Its Price}
    {Thomas~Wein}
    {\textit{Journal of Industrial Economics}, 2025, 73(2): 351--383.}
  \item ``\href{http://www.zeitschrift-fuer-verkehrswissenschaft.de/fileadmin/archiv/hefte---2025_1_2_3/2025-1/02a_ZfV_Heft_2025-1_Beitrag_Kahl_Wein.pdf}{Existiert Wettbewerb beim Laden von Elektrofahrzeugen in Deutschland?}''
    [Is There Competition for Charging Electric Vehicles in Germany?] (with~Thomas~Wein)\\
    \textit{Zeitschrift für Verkehrswissenschaft}, 2025, issue~1: 31--66. In German.
  \paper{https://doi.org/10.1016/j.eneco.2024.107961}
    {Cross-Border Competition in the Gasoline Retail Market: Impact of Proximity at the German--Polish Border}
    {}
    {\textit{Energy Economics}, 2024, 140: 107961.}
  \paper{https://doi.org/10.1016/j.eneco.2024.107843}
    {Was the German Fuel Discount Passed On to Consumers?}
    {}
    {\textit{Energy Economics}, 2024, 138: 107843.}
\end{papers}

% ---- Working papers ----------------------------------------------------
\section{Working Papers}
\begin{papers}
  \paper{https://www.leuphana.de/fileadmin/user_upload/Forschungseinrichtungen/ifvwl/WorkingPapers/lue/pdf/wp_446_Upload.pdf}
    {Three Clocks on a Tax Date: Cost, Posting, and the Measurement of Tax Pass-Through}
    {Thomas~Wein}
    {\textit{University of Lüneburg Working Paper Series in Economics}, No.~446, 2026.}
  \paper{https://www.leuphana.de/fileadmin/user_upload/Forschungseinrichtungen/ifvwl/WorkingPapers/wp_444_Upload.pdf}
    {More Competition (Only) Through Greater Market Transparency? Twelve Years of the Fuel Market Transparency Agency in Germany}
    {Thomas~Wein}
    {\textit{University of Lüneburg Working Paper Series in Economics}, No.~444, 2026.}
\end{papers}

% ---- Work in progress --------------------------------------------------
\section{Work in Progress}
\begin{papers}
  \wip{Strategic Pricing and Consumer Welfare under One-Sided Price Regulation: Theory and Evidence from German Retail Fuel Markets}{Philipp~Denter}
  \wip{AI and the Endogenous Organization of Social Interactions at Work: A Theoretical and Empirical Investigation}{Facundo~Albornoz, Antonio~Cabrales, and Esther~Hauk}
  \wip{Electric Vehicle Subsidy Pass-Through: Evidence from Germany}{Adrián~Santonja~di~Fonzo}
  \wip{Third-Party Effects of Subsidies: The Case of German Subsidies for Cafeteria Food}{Jonas~Dix}
  \wip{Does Import Origin Determine Who Pays? Import Exposure and the Cross-Country Incidence of Europe's 2022 and 2026 Diesel Shocks}{}
  \wip{Weather Meets Geopolitics: Rhine Low Water, Inland Fuel Supply, and the Regional Incidence of an Oil Shock}{}
  \wip{Gas Lines Are Back: A Voluntary Price Cap, Stockouts, and Non-Price Rationing in the 2026 French Fuel Crisis}{}
\end{papers}

% ---- Presentations -----------------------------------------------------
\section{Conference and Seminar Presentations}
\begin{entries}
  \entry{2026}{Annual Congress of the European Economic Association (Dublin);
    Applied Reading Group, Universidad Carlos III de Madrid;
    Research Seminar, Universidad Autónoma de Madrid;
    Simposio de la Asociación Española de Economía (Las Palmas de Gran Canaria, accepted)}
  \entry{2024}{EARIE Annual Conference (Amsterdam, accepted);
    6th Workshop for Young Scientists on Applied Economics (Hannover)}
  \entry{2023}{EARIE Annual Conference (Rome);
    5th Workshop for Young Scientists on Applied Economics (Hannover);
    Doctoral Program Behavioral Economics and Societal Transformation (Lüneburg);
    Empirical Microeconomics PhD Workshop (Lüneburg)}
  \entry{2022}{Empirical Microeconomics PhD Workshop (Lüneburg)}
  \entry{2021}{Annual Meeting of the Verein für Socialpolitik;
    Empirical Microeconomics PhD Workshop (Lüneburg)}
\end{entries}

% ---- Awards ------------------------------------------------------------
\section{Honors and Scholarships}
\begin{entries}
  \entry{2025}{Leuphana Dissertation Award, Leuphana University Lüneburg}
  \entry{2017--2024}{e-fellows.net Scholarship}
\end{entries}

% ---- Teaching ----------------------------------------------------------
\section{Teaching Experience}
\begin{entries}
  \entry{2025--present}{\textbf{Universidad Carlos III de Madrid}, Lecturer (in English)\\
    Economics I (Microeconomics); Industrial Organization}
  \entry{2022--present}{\textbf{Verwaltungs- und Wirtschaftsakademie (VWA) and Berufsakademie (BA) Lüneburg},
    Lecturer (in German)\\
    Macroeconomics}
  \entry{2019--2025}{\textbf{Leuphana University Lüneburg}\\
    \textit{Graduate School, own courses (M.A., in English):} Introduction to Economics for
      Non-Economists; Economics Preparatory Course\\
    \textit{Professional School, junior lecturer (in English):} Regulation Economics; Competition Economics\\
    \textit{College, teaching assistant (B.Sc., in English and German):} Macroeconomics~1 and~2;
      Microeconomics~2; Applied Macroeconomics: International Economics;
      Applied Microeconomics: Industrial Economics}
  \entry{2016--2018}{\textbf{Leuphana University Lüneburg}, Tutor in Microeconomics (in English),
    Faculty of Economics and Center for the Study of Democracy}
\end{entries}

% ---- Service -----------------------------------------------------------
\section{Refereeing}
\begin{entries}
  \entry{}{Annals of Regional Science; Economics Letters; Energy Economics; The Energy Journal;
    Finance Research Letters; Journal of Baltic Studies}
\end{entries}

\section{University Service}
\begin{entries}
  \entry{2019--2025}{\textbf{Leuphana University Lüneburg}\\
    Program Coordinator, Economics Minor (incl.\ reaccreditation, 2023)\\
    Deputy Program Coordinator, B.Sc.\ Economics (incl.\ conversion into an
      English-language program, 2022, and reaccreditation, 2023)\\
    Member, Examination Board\\
    Member, Study Commission Governance \& Law (2019--2022)\\
    Member, Appointment Committees for the Professorships in Law and Economics (2023) and in
      Applied Economic Policy/Macroeconomics (2022)}
  \entry{2022--present}{\textbf{Alumni- und Förderverein der Leuphana Universität Lüneburg e.V.}\
    (alumni association), Board Member; Chair since 2024}
\end{entries}

% ---- Memberships -------------------------------------------------------
\section{Memberships}
\begin{entries}
  \entry{2024--present}{European Economic Association}
  \entry{2021--present}{Verein für Socialpolitik (German Economic Association)}
\end{entries}

% ---- Training ----------------------------------------------------------
\section{Additional Training}
\begin{entries}
  \entry{2025}{BSE Summer School, Barcelona School of Economics: ``Dynamic and Non-linear Panel
    Data Models'' (Badi Baltagi); ``Digital Empirical Methods for Cartel and Corruption Screening
    in Procurement'' (Kei Kawai and Juan José Ganuza)}
  \entry{2025}{PhD course ``Causal Inference and Applied Microeconometrics'' (Nils Braakmann),
    Leuphana University Lüneburg}
  \entry{2024}{BSE Summer School: ``Difference-in-Differences with Panel Data'' (Jeffrey M.\ Wooldridge)}
  \entry{2024}{PhD course ``Field Experiments'' (Adriaan Soetevent), University of Vechta}
  \entry{2023}{BSE Summer School: ``Empirical Methods for the Analysis of the Energy Transition''
    (Mar Reguant)}
  \entry{2022}{BSE Summer School: ``Quantitative Methods for Public Policy Evaluation'' (Albrecht
    Glitz); ``Dynamic and Non-linear Panel Data Models'' (Sergi Jiménez-Martín)}
  \entry{2021}{ÊSTIMATE Summer School (Jeffrey M.\ Wooldridge), Michigan State University}
  \entry{2020}{14th International Research Workshop, Helmut Schmidt University Hamburg, incl.\
    ``Spatial Data Analysis with Stata'' (Timo Mitze)}
\end{entries}

% ---- Skills ------------------------------------------------------------
\section{Skills}
\begin{entries}
  \entry{Software}{Stata, R, LaTeX, Python}
  \entry{Languages}{German (native), English (fluent), Spanish (basic), French (beginner)}
\end{entries}

% ---- References --------------------------------------------------------
\section{References}
\begin{papers}
  \item
  \referee{Antonio Cabrales}{Professor of Economics}
    {Department of Economics\\Universidad Carlos III de Madrid}{antonio.cabrales@uc3m.es}\hspace{0.04\linewidth}%
  \referee{Philipp Denter}{Ramón y Cajal Researcher}
    {Department of Economics\\Universidad Carlos III de Madrid}{pdenter@eco.uc3m.es}

  \vspace{1.6ex}
  \referee{Boris Hirsch (second Ph.D.\ advisor)}{Professor of Economics}
    {Institute of Economics\\Leuphana University Lüneburg}{boris.hirsch@leuphana.de}\hspace{0.04\linewidth}%
  \referee{Thomas Wein (Ph.D.\ advisor)}{Professor of Economics}
    {Institute of Economics\\Leuphana University Lüneburg}{thomas.wein@leuphana.de}
\end{papers}

\end{document}
